\section*{CHAPTER 5: TESTING, DEPLOYMENT, AND CONCLUSION}
\addcontentsline{toc}{section}{\numberline{}CHAPTER 5: TESTING, DEPLOYMENT, AND CONCLUSION}
\setcounter{section}{5}
\setcounter{subsection}{0}
\setcounter{figure}{0}
\setcounter{table}{0}

This final chapter outlines the quality assurance processes, the cloud infrastructure used to deploy the platform, and provides a comprehensive evaluation of the project's successes, limitations, and future potential.

\subsection{System Testing}
Ensuring the reliability of a financial and scheduling system is paramount. The testing phase was divided into three main categories:
\begin{itemize}
    \item \textbf{Unit Testing:} Individual functions, particularly the occupancy calculation algorithms and multi-tenant scoping logic, were isolated and tested to ensure mathematical accuracy and data security.
    \item \textbf{Integration Testing:} The interactions between the Node.js backend and the MySQL database were tested to verify that the Knex.js queries correctly handled edge cases, such as attempting to book a field that overlaps with an existing booking.
    \item \textbf{User Acceptance Testing (UAT):} The mobile application and web dashboards were tested on various physical devices (iOS, Android, Windows) to ensure responsive design compliance and verify that the QR code scanner operated smoothly under different lighting conditions.
\end{itemize}

\subsection{Cloud Deployment Architecture}
The FootField platform is designed to be fully cloud-native, ensuring high availability and eliminating the need for on-premises server maintenance.

The \textbf{Node.js Application} (encompassing both backend APIs and frontend static files) is deployed on \textbf{Render}, a modern Platform-as-a-Service (PaaS). A Continuous Integration/Continuous Deployment (CI/CD) pipeline is established by linking Render directly to the project's GitHub repository. Any code pushed to the \texttt{main} branch automatically triggers a build and deployment process on Render's servers, ensuring the production environment is always up-to-date with the latest verified code.

The \textbf{MySQL Database} is decoupled from the application server and hosted on \textbf{Aiven}, a specialized Database-as-a-Service (DBaaS) provider. This separation ensures that database compute resources (CPU/RAM) are not competing with the web server. Aiven also provides automated daily backups and point-in-time recovery, safeguarding crucial financial and booking data against accidental loss.

\subsection{Evaluation of Results}
The deployed platform successfully meets the core objectives defined at the project's inception. 
\begin{enumerate}
    \item The \textbf{Multi-Tenant Architecture} allows the system to seamlessly host multiple sports facilities on a single codebase and database, proving the economic viability of the SaaS model.
    \item The \textbf{AI Chatbot} significantly reduces the manual workload of facility managers. By allowing customers to check availability and book fields autonomously via natural language, the system operates effectively 24/7.
    \item The \textbf{Capacitor Mobile App} successfully bridges the software-hardware gap, transforming standard smartphones into dedicated QR scanning terminals for rapid on-site check-ins.
\end{enumerate}

\subsection{Limitations}
Despite its successes, the current iteration exhibits certain limitations:
\begin{itemize}
    \item \textbf{AI Latency:} Because the AI chatbot relies on the external Google Gemini API, response times can occasionally exceed 2 seconds during peak network congestion, which may momentarily disrupt the conversational flow.
    \item \textbf{Payment Gateway Exclusivity:} The system currently only integrates VNPay. While popular, the lack of support for other major e-wallets (e.g., MoMo, ZaloPay) or international credit cards limits user payment options.
\end{itemize}

\subsection{Future Work}
To evolve FootField into a comprehensive enterprise solution, several future development paths have been identified:
\begin{itemize}
    \item \textbf{IoT Smart Lighting Integration:} Developing hardware relays connected to the booking database. This would allow the system to automatically turn on the football field floodlights exactly when a booked session begins and turn them off when it ends, drastically reducing electricity waste.
    \item \textbf{Machine Learning Analytics:} Implementing predictive models to analyze historical booking data. This would enable the system to automatically suggest dynamic pricing strategies (e.g., increasing prices during predicted high-demand hours) to maximize revenue for facility owners.
\end{itemize}

\subsection{Conclusion}
In conclusion, the "FootField" project demonstrates the transformative potential of combining modern web architectures, Artificial Intelligence, and mobile integration within traditional service industries. By replacing fragmented, manual management processes with a centralized, automated SaaS platform, the project successfully bridges the gap between technology and sports facility operations. The resulting system is not only a robust academic exercise but also a highly viable commercial product ready for real-world deployment.
